% ============================================================
% Chapter 2 — Background (target 8–10 pp)
% Status: FULL DRAFT 2026-07-14. All 12 bibliography entries verified
% against publisher pages on 2026-07-14 (outline open item 5 closed).
% Scope rule: only concepts actually used in Ch. 3–6, each with a
% forward reference to where it is used.
% ============================================================
\chapter{Background}
\label{ch:background}

This chapter introduces the concepts the rest of the thesis builds on:
single-agent reinforcement learning and policy gradients
(\cref{sec:rl-essentials}), PPO and advantage estimation
(\cref{sec:ppo}), multi-agent learning and MAPPO (\cref{sec:marl}),
curriculum learning and its mixed-training control
(\cref{sec:curriculum-bg}), and the CARLA simulator
(\cref{sec:carla-bg}). Each concept is presented only to the depth at
which it is used later.

% ------------------------------------------------------------
\section{Reinforcement Learning Essentials}
\label{sec:rl-essentials}

Reinforcement learning (RL) formalizes sequential decision making as a
\emph{Markov decision process} (MDP): a tuple
$(\mathcal{S}, \mathcal{A}, P, r, \gamma)$ of states, actions, transition
dynamics $P(s' \mid s, a)$, reward function $r(s, a)$, and discount
factor $\gamma \in [0, 1)$~\cite{sutton2018rl}. An agent acts through a
\emph{policy} $\pi(a \mid s)$, and the learning objective is to maximize
the expected discounted return
\begin{equation}
  J(\pi) \;=\; \mathbb{E}_{\pi}\!\left[\sum_{t=0}^{\infty} \gamma^{t}\,
  r(s_t, a_t)\right].
\end{equation}
The discount sets an effective decision horizon of roughly
$1/(1-\gamma)$ steps. This is not a cosmetic constant: the platform of
\cref{ch:system} uses $\gamma = 0.997$, i.e.\ an effective horizon of
$\approx 333$ steps, commensurate with episodes of up to 1{,}000 control
steps in which sparse route-progress rewards may be tens of seconds
apart.

Two value functions summarize a policy: the state value
$V^{\pi}(s) = \mathbb{E}_{\pi}[\,\text{return} \mid s_0 = s\,]$ and the
action value $Q^{\pi}(s, a)$; their difference
$A^{\pi}(s, a) = Q^{\pi}(s, a) - V^{\pi}(s)$, the \emph{advantage},
measures how much better an action is than the policy's average in that
state.

\paragraph{Policy gradients and actor--critic methods.}
When the policy is a differentiable function $\pi_{\theta}$ (here, a
neural network), the policy-gradient theorem gives an unbiased gradient
estimate of the objective:
\begin{equation}
  \nabla_{\theta} J \;=\;
  \mathbb{E}_{\pi_{\theta}}\!\left[\nabla_{\theta}
  \log \pi_{\theta}(a_t \mid s_t)\; A^{\pi}(s_t, a_t)\right].
\end{equation}
Estimating the advantage with raw Monte-Carlo returns yields high
variance; subtracting a learned state-value baseline reduces variance
without introducing bias. \emph{Actor--critic} methods therefore train
two functions: the \emph{actor} $\pi_{\theta}$ and a \emph{critic}
$V_{\phi}$ used to construct advantage estimates. This division of labour
is the hook for the multi-agent architecture of \cref{sec:marl}: the
critic is a training-time device, so it may receive information the
actor never sees.

\paragraph{Continuous actions.}
For continuous control the policy is commonly a diagonal Gaussian: the
network outputs a mean (and log-standard-deviation) per action dimension,
and actions are sampled around it. Sampling is not an implementation
detail but part of the learned object --- evaluating such a policy at its
deterministic mean discards its stochastic structure, a point that
becomes concrete in \cref{sec:bugs}.

% ------------------------------------------------------------
\section{Proximal Policy Optimization}
\label{sec:ppo}

Vanilla policy-gradient updates are fragile: a single large step can
collapse a working policy, and on-policy data is stale immediately after
the update. PPO~\cite{schulman2017ppo} addresses both with a
\emph{clipped surrogate objective}. With the importance ratio
$\rho_t(\theta) = \pi_{\theta}(a_t \mid s_t) / \pi_{\theta_{\text{old}}}(a_t \mid s_t)$,
PPO maximizes
\begin{equation}
  L^{\mathrm{CLIP}}(\theta) \;=\;
  \mathbb{E}_t\!\left[\min\bigl(\rho_t A_t,\;
  \operatorname{clip}(\rho_t,\, 1-\varepsilon,\, 1+\varepsilon)\, A_t\bigr)\right],
\end{equation}
which removes any incentive to move the ratio beyond
$[1-\varepsilon, 1+\varepsilon]$, allowing several epochs of minibatch
gradient steps on each batch of experience without destructive updates.
Implementations often add an adaptive Kullback--Leibler penalty toward
the pre-update policy as a second safeguard; the configuration in
\cref{sec:optimization} uses the clip ($\varepsilon = 0.2$) and a KL
target of $0.02$ together.

\paragraph{Generalized Advantage Estimation.}
The advantage estimates $A_t$ come from GAE~\cite{schulman2016gae}:
with temporal-difference residuals
$\delta_t = r_t + \gamma V(s_{t+1}) - V(s_t)$,
\begin{equation}
  A_t^{\mathrm{GAE}(\gamma,\lambda)} \;=\;
  \sum_{l \geq 0} (\gamma\lambda)^{l}\, \delta_{t+l},
\end{equation}
where $\lambda \in [0,1]$ trades bias against variance between one-step
TD ($\lambda = 0$) and Monte-Carlo returns ($\lambda = 1$). The platform
uses $\lambda = 0.95$. GAE quality depends directly on the critic's
quality --- one documented development episode
(\cref{sec:iterations}) involved a critic whose explained variance was
$\approx 0$, making advantages nearly baseline-free until the
value-loss configuration was repaired.

\paragraph{Entropy regularization.}
An entropy bonus added to the objective sustains exploration by
penalizing premature policy collapse. Its coefficient is typically
annealed: high early (broad exploration), low late (exploitation and
stable convergence). In this thesis the coefficient follows an explicit
schedule from $0.03$ to $0.005$ anchored at a fixed \emph{fraction} of
the training budget (\cref{sec:optimization}), so that runs of different
lengths traverse the same exploration profile.

% ------------------------------------------------------------
\section{Multi-Agent Reinforcement Learning and MAPPO}
\label{sec:marl}

With $N$ agents acting in one environment, the MDP generalizes to a
\emph{Markov game}: transitions and rewards depend on the joint action.
Two structural difficulties follow.

\paragraph{Non-stationarity.}
From any single agent's perspective, the other agents are part of the
environment --- and they are learning, so the environment's effective
dynamics change over time. This breaks the stationarity assumptions
underlying single-agent convergence guarantees and makes
experience collected under old co-player policies
misleading~\cite{lowe2017maddpg}. In the platform of \cref{ch:system}
this is not a corner case but the core of the setting: six learning
agents share one world, and \cref{ch:results} documents a concrete
equilibrium (vehicles freezing in response to pedestrian behaviour) that
only exists because of this coupling.

\paragraph{Credit assignment and variance.}
Each agent's return depends on all agents' actions, so the variance of
naive policy-gradient estimates grows with the number of agents: an
agent may be rewarded or punished for outcomes it did not cause.

\paragraph{Centralized training, decentralized execution.}
The now-standard resolution exploits the actor--critic division: train
\emph{critics} with access to global information (all agents'
observations, or the true state), while keeping \emph{actors}
conditioned only on local observations~\cite{lowe2017maddpg}. The critic
absorbs the non-stationarity --- conditioned on everyone's situation, the
value target is better determined --- while execution remains fully
decentralized, since critics are discarded after training. This is the
CTDE paradigm used throughout this thesis.

\paragraph{MAPPO.}
MAPPO instantiates CTDE with PPO actors and a centralized value
function. Yu et al.~\cite{yu2022surprising} showed that this
conceptually minimal combination --- long assumed too sample-inefficient
for MARL --- matches or exceeds off-policy methods on standard
cooperative benchmarks (particle worlds, StarCraft micromanagement,
Hanabi, football), provided the implementation details are right; among
the factors they highlight are the design of the centralized critic's
input and value normalization. Those two factors reappear concretely in
this thesis as the global-observation layout of \cref{sec:obs-actions}
and the value-loss configuration episode of \cref{sec:iterations}.

\paragraph{Parameter sharing and heterogeneous classes.}
Homogeneous agents can share one policy network, pooling their
experience; heterogeneous agents cannot, since their observation and
action spaces differ. The platform takes the intermediate route natural
to its population: one shared policy per class (three vehicles; three
pedestrians), two policies in total (\cref{sec:architecture}).

\paragraph{Critic input encoders.}
The simplest centralized critic flattens all agents' observations into
one vector for an MLP. Structured alternatives treat agents as nodes of
a set or graph: per-agent slots are embedded, combined by attention or
by message passing, and pooled. The graph-based variant used as this
thesis's secondary experimental axis applies GraphSAGE-style mean
aggregation~\cite{hamilton2017graphsage} over the six agent nodes with
an alive mask (\cref{sec:critic-encoders}). The comparison MLP-vs-GNN
critic asks whether relational structure in value estimation changes
what the (identical) actors learn.

\paragraph{Coexistence rather than cooperation.}
Unlike the benchmarks above, the agents of this thesis neither share a
team reward nor compete: each pursues its own route under an individual
reward (\cref{sec:rewards}). CTDE remains the right tool --- the
centralized critic still stabilizes value estimation under mutual
interaction --- but the setting is best described as \emph{coexistence}
of independent tasks in a shared world, the framing formalized as
parallel task coverage in \cref{sec:parallel-coverage}.

% ------------------------------------------------------------
\section{Curriculum Learning and Its Control Condition}
\label{sec:curriculum-bg}

\paragraph{Origins.}
The idea that learning improves when examples are presented in a
meaningful easy-to-hard order predates deep learning: Elman's
``starting small'' experiments showed recurrent networks learning a
complex grammar only when input complexity (or the network's effective
memory) grew incrementally~\cite{elman1993starting}. Bengio et
al.~\cite{bengio2009curriculum} named the strategy \emph{curriculum
learning} and framed it as a continuation method: easy examples define a
smoother version of the objective on which optimization starts, hard
examples are introduced as training proceeds, with measurable gains in
optimization speed and generalization on supervised tasks.

\paragraph{Curriculum learning in RL.}
In reinforcement learning the designer controls not just the order of
data but the distribution of \emph{tasks}, making curricula both more
powerful and harder to design. Narvekar et al.~\cite{narvekar2020curriculum}
organize the field around three questions: how source tasks are
generated, how they are \emph{sequenced}, and how knowledge transfers
between them. Within that taxonomy, the mechanism of \cref{sec:regimes}
is a performance-based sequencing scheme over a fixed task set: three
difficulty levels ordered by route length, with progression gated on a
windowed success threshold and guaranteed by budget-based pressure caps.
The motivation for curricula in RL is largely about exploration: on hard
tasks a weak policy may essentially never experience success, receiving
no learning signal at all, while easy tasks provide dense early success
from which competence can bootstrap.

\paragraph{Known risks.}
Three failure modes recur in the literature and shaped the design of
\cref{sec:regimes}. First, \emph{forgetting}: neural networks trained on
a shifted distribution tend to overwrite previously acquired
competence~\cite{french1999catastrophic}; a curriculum that abandons
easy tasks after progression may erode exactly the skills it built ---
the platform therefore keeps all unlocked levels in the sampling mix
rather than replacing them. Second, \emph{stalling}: a competence gate
that is never met can trap training on easy tasks indefinitely ---
hence the pressure caps that force progression within the budget.
Third, \emph{gate mis-calibration}: unlocking too early exposes the
policy to noise it cannot yet learn from; unlocking too late wastes
budget. These risks are exactly the hypotheses H1 tests empirically
(\cref{sec:rq}).

\paragraph{The control condition: mixed training.}
Every claim about ordering needs a baseline in which ordering is absent.
The natural control is \emph{mixed} (interleaved) training: all
difficulty levels presented from the first episode, with uniform
exposure --- the RL analogue of training on a shuffled dataset, which is
the default in supervised learning. The stratified sampler of
\cref{sec:batch-arm} implements this control exactly (uniform $1/3$
exposure with maximum imbalance of one episode), so that the
curriculum-versus-mixed comparison isolates the ordering variable at
equal budget. Mixed training carries its own hypothesized signature
(H2): full task coverage from the start, at the price of a noisier early
signal dominated by failures on hard tasks.

% ------------------------------------------------------------
\section{The CARLA Simulator}
\label{sec:carla-bg}

CARLA~\cite{dosovitskiy2017carla} is an open-source urban-driving
simulator built on Unreal Engine~4, developed to support training and
validation of autonomous driving systems. It runs as a server hosting
the 3D world and physics, controlled by clients through a Python API;
it ships with a set of towns of varying layout, scripted traffic, and
configurable weather. Three features matter for this thesis
(version 0.9.16 throughout):

\begin{itemize}
  \item \textbf{Synchronous mode with fixed time step.} The server
        advances only on client ticks, at a fixed physics step
        (\SI{0.05}{s} here). Experience collection is thus reproducible
        in structure and independent of rendering or wall-clock speed
        (\cref{sec:architecture}).
  \item \textbf{Traffic Manager and walker control.} Scripted (NPC)
        vehicles are driven by CARLA's Traffic Manager; pedestrians are
        \emph{walker} actors that accept direct control commands, which
        is what makes learning pedestrian policies possible
        (\cref{sec:obs-actions}).
  \item \textbf{Structurally different towns.} Training uses
        \emph{Town03}, described by the official documentation as ``a
        larger, urban map with a roundabout and large
        junctions''~\cite{carladocs}; the generalization test uses
        \emph{Town05}, a ``squared-grid town with cross junctions and a
        bridge'', with multiple lanes per direction~\cite{carladocs}.
        The layouts differ in kind, not just in size, making the
        Town03$\to$Town05 transfer a genuine test of map generalization
        (\cref{sec:generalization}).
\end{itemize}

Two scoping notes. First, in this thesis CARLA provides physics, world
dynamics and traffic --- not perception: agents observe privileged state
features (\cref{sec:obs-actions}), not camera or lidar streams, so all
findings concern decision-making under interaction, not vision. Second,
even in synchronous mode and fully seeded, CARLA under multi-agent load
is not run-to-run deterministic; this irreducible variability is
quantified with same-seed replicates and used as the yardstick for every
effect size reported later (\cref{sec:iterations},
\cref{sec:threats}).
